\section{Low-Precision Output Head: Negative Result}
\label{app:output-head}

This appendix records an engineering screen for which the exact low-precision
forward format was not preserved. The vocabulary projection has a large
forward GEMM and is therefore an attractive FP4 target. In the complete
trainer, however, the head also owns cross entropy, the hidden-state gradient,
the weight gradient, quantization, and saved-tensor lifetime. Accelerating one
GEMM did not accelerate that whole boundary.

MXFP8 below denotes an eight-bit microscaling format used for one
weight-gradient configuration.

\begin{table}[H]
  \centering
  \caption{Historical whole-trainer output-head screen. Deltas are paired
  within one experiment. The complete profiler metadata was not preserved, so
  rates are rounded and are not compared with the primary performance panel.}
  \label{tab:head-negative}
  \begin{tabularx}{\textwidth}{YrrY}
    \toprule
    Forward / hidden-state gradient / weight gradient & Tokens/s/GPU & Difference & Outcome \\
    \midrule
    BF16 / BF16 / BF16 & $\sim$34.00K & baseline & BF16 control \\
    lowp$^{\dagger}$ / BF16 / BF16 & $\sim$33.95K & $\sim-0.16\%$ & Loss differed from control \\
    lowp$^{\dagger}$ / BF16 / MXFP8 & $\sim$33.91K & $\sim-0.27\%$ & Loss and gradients differed \\
    \bottomrule
  \end{tabularx}
\end{table}

$^{\dagger}$ The experiment record preserves the backward choices and paired
rates but not the exact low-precision forward format. We therefore retain the
literal ``lowp'' label rather than assigning a more specific format after the
fact.

Both selective low-precision variants are slower than their matched BF16
control. The preserved record also says that neither reproduced the BF16
control loss, and that the MXFP8 weight-gradient configuration additionally
changed gradients; it
does not retain a publication-ready error magnitude, so we do not rank the
size of those numerical differences. These results support the selected
design: a BF16 output projection with compiled CE.
